



\chapter{Promise Body Parts}





\section{Body Parts}

A promise body is the description of exactly what is promised (as
opposed to what/who is making the promise). The term \textbf{body} (or
\textbf{body part}) is used in the CFEngine syntax to mean a small template
that can be used to contribute as part of a larger promise body.

\begin{description}
\item[\textbf{Body}] A body part is a collection of promise attributes.
\end{description}

\begin{codelisting}
\codecaption{400-010-Body\_Parts-0010-Introduction.cf}
%= lang:cfengine3, options: "linenos": false
\begin{code}
body <type> name
{
  attribute1 => value1;
  attribute2 => value2;
  ...
  attributeN => valueN;
}
\end{code}
\end{codelisting}
\begin{codelisting}
\codecaption{400-010-Body\_Parts-0500-Parts\_No\_world\_write\_bit.cf}
%= lang:cfengine3, options: "linenos": false
\begin{code}
# run "groupadd project_team" to create the group
# "project_team" for this example

bundle agent main
{
  files:

      "/tmp/file1"
        create => "true",
        perms   => project_files;

      "/tmp/file2"
        create => "true",
        perms   => project_files;
}

body perms project_files
{
        mode   => "770";
        owners => { "root" };
        groups => { "project_team" };
}
\end{code}
\end{codelisting}
\begin{codelisting}
\codecaption{400-010-Body\_Parts-0540-Parts\_Setting\_group\_ownership\_based\_on\_OS.cf}
%= lang:cfengine3, options: "linenos": false
\begin{code}
# You can use class expressions in body parts

bundle agent main
{
  files:
    any::
      "/tmp/testfile"
        comment  => "Set appropriate file attributes for Admin group
                     on every type of OS",
        create  => "true",
        perms   => admin_group;
}

body perms admin_group
{
  linux::  groups => { "wheel" };
  darwin:: groups => { "admin" };
  sunos::  groups => { "sys" };

}
\end{code}
\end{codelisting}
\begin{codelisting}
\codecaption{400-010-Body\_Parts-0545-parameterized\_body.cf}
%= lang:cfengine3, options: "linenos": false
\begin{code}
# You can parameterize body parts

bundle agent main
{
  files:

      "/tmp/testfile"
        create  => "true",
        perms   => set_mode_700_admin_group_and_specified_user("sam");
}

body perms set_mode_700_admin_group_and_specified_user(user)
{
        mode   => "0700";
        owners => { "$(user)" };

      linux::  groups => { "wheel" };
      darwin:: groups => { "admin" };
      sunos::  groups => { "sys" };
}
\end{code}
\end{codelisting}
\begin{codelisting}
\codecaption{400-010-Body\_Parts-0550-Parts\_Replacing\_Patterns.cf}
%= lang:cfengine3, options: "linenos": false
\begin{code}
# Run "echo I am going to walk my dog > /tmp/data.txt" to set up for
# this example

# Demonstrate search-and-replace in a file

bundle agent main
{
  files:
    any::
      "/tmp/data.txt"
        edit_line => transform_dogs_to_cats;
}

bundle edit_line transform_dogs_to_cats
{
  replace_patterns:
      "[Dd]og"
        replace_with => value("cat");
}

body replace_with value(x)
{
        replace_value => "$(x)";
        occurrences => "all";
}
\end{code}
\end{codelisting}
\begin{codelisting}
\codecaption{400-010-Body\_Parts-0580-perms.cf}
%= lang:cfengine3, options: "linenos": false
\begin{code}
# files: perms takes a list of okay groups
# Change to first entry if not matching any of the groups in the list
#
# Run "sudo touch /tmp/testfile; sudo chmod 777 /tmp/testfile; sudo chgrp nobody /tmp/testfile"
# to set up for this example.

bundle agent main
{
  files:
      "/tmp/testfile"
        perms => not_world_writable_and_right_group;
}

body perms not_world_writable_and_right_group
{
        groups => {"root", "games", "mail" };
        mode   => "o-w";
}
\end{code}
\end{codelisting}
\begin{codelisting}
\codecaption{400-010-Body\_Parts-0590-perms.cf}
%= lang:cfengine3, options: "linenos": false
\begin{code}
# A body part can be re-used across promises
#
# run "useradd bob; useradd susan" to set up for this example

bundle agent main
{
  files:
      "/tmp/bobsfile"
        create  => "true",
        perms   => mog("777", "bob", "mail");

      "/tmp/susansfile"
        create  => "true",
        perms   => mog("000", "susan", "games");
}

body perms mog(mode,owner,group)
{
  mode    => "$(mode)";
  owners  =>  {"$(owner)"};
  groups  =>  {"$(group)"};
}
\end{code}
\end{codelisting}





\begin{aside}
\label{aside:exercise_29}
\heading{Create executable shell script}

Write a CFEngine policy to ensure `/usr/local/bin/hello.sh' exists,
has permissions 0755, owner root, group root, and contents:

%= lang:bash
\begin{code}
#!/bin/sh
/bin/echo hello world
\end{code}

\end{aside}





\section{Left Hand Side, Right Hand Side}

\begin{codelisting}
\codecaption{400-020-CFEngine\_Grammar-0100-Fix\_Your\_Webserver.cf}
%= lang:cfengine3, options: "linenos": false
\begin{code}
# Example of
#     cfengine_word => builtin_function()

bundle agent main
{
  vars:
      "formatted"
        string => format("%04d", 1);

  reports:
      "$(formatted)";
}
\end{code}
\end{codelisting}
\begin{codelisting}
\codecaption{400-020-CFEngine\_Grammar-0120-List\_Example\_B\_Var\_Slist\_Only.cf}
%= lang:cfengine3, options: "linenos": false
\begin{code}
# Example of
#
#      cfengine_word => { list }    #  (directly and via list variable)

bundle agent main
{
  vars:
      "my_slist_0"
        slist  => {
                    "String contents...",
                    "... are beauutifuuul this time of year"
                  };

      "my_slist_1"
        slist  => { @(my_slist_0), "apple", "orange" };
}
\end{code}
\end{codelisting}





\section{LHS/RHS (Continued)}





\subsection{Promise attributes}

CFEngine uses many ``constraint expressions'' as part of the body of a
promise.  These are attributes of a promise, they detail and constrain
the promise.

These take the form:

left-hand-side (cfengine word) =\textgreater{} right-hand-side (user defined data).

This can take several forms:

\begin{verbatim}
 cfengine_word => user_defined_body or user_defined_body(parameters)

                  builtin_function()

                  "scalar_value" or "$(scalar_variable_name)"

                  { "list_element", "list_element2" }

                  { @(list_variable_name) }

                  boolean
\end{verbatim}

In each of these cases, the right hand side is a user choice.

\begin{codelisting}
\codecaption{400-020-CFEngine\_Grammar-0160-LHS\_vs\_RHS\_Example\_of\_user\_defined\_body\_on\_rhs\_WITH\_PARAMS.cf}
%= lang:cfengine3, options: "linenos": false
\begin{code}
# Example of:
#    cfengine_word => user_defined_body(param)

bundle agent main
{
  storage:
      "/" volume  => my_check_volume("30%", "100K");
      "/var" volume  => my_check_volume("20%", "500K");
}

body volume my_check_volume(min_free_space,size)
{
        freespace      => "$(min_free_space)";
        # Min disk space that should be available

        sensible_size  => "$(size)";
        # Minimum size in bytes that should be used
}
\end{code}
\end{codelisting}
\begin{codelisting}
\codecaption{400-020-CFEngine\_Grammar-0180-LHS\_vs\_RHS\_Example\_of\_scalar\_on\_rhs.cf}
%= lang:cfengine3, options: "linenos": false
\begin{code}
# Example of
#      cfengine_word => "scalar"

bundle agent main
{
  vars:

      "variable_0"
        comment => "RHS is a literal string",
        string  => "String contents...";  # a scalar value

      "variable_1"
        comment => "RHS uses a variable but it's still a scalar",
        string  => "$(variable_0)";       # a scalar variable

  reports:
      "variable_0: $(variable_0)";
      "variable_1: $(variable_1)";
}
\end{code}
\end{codelisting}
\begin{codelisting}
\codecaption{400-020-CFEngine\_Grammar-0190-LHS\_vs\_RHS\_Example\_of\_list\_on\_rhs.cf}
%= lang:cfengine3, options: "linenos": false
\begin{code}
# Example of
#
#      cfengine_word => { list }    #  (directly and via variable)

body common control
# "body common control" contains attributes that control the behavior
# of CFEngine
{

  bundlesequence => { "example_1", "example_2" };

}

bundle agent example_2
{
  reports:
      "Second things second";
}

bundle agent example_1
{
  reports:
      "First things first";
}
\end{code}
\end{codelisting}